\section{Ligninger}

\subsubsection*{Eksempel 1)}
Der er mange måder at skrive ligninger! Blandet andet kan du bruge \textbf{\textbackslash begin \{equation\}}:
\begin{equation}
    y=5x+\frac{1}{2}
\end{equation}

\begin{equation*}
    y=5x+\frac{1}{2}
\end{equation*}

$$y=5x+\frac{1}{2}$$

\[y=5x+\frac{1}{2}\]

Alternativt kan du bruge \textbf{\textbackslash \{align\}} for at få ligningerne lige under hinanden:
\begin{align*}
    x&=12\\
    x131212&=12213
\end{align*}

\begin{align}
    x&=12\\
    x131212&=12213
\end{align}

\subsubsection*{Eksempel 2)}
Du kan også tilgå "matematik-mode" ved at bruge dollar tegn omkring matematikken \textbf{\$ \$}: \\
Normal Approximation to the Binomial Distribution
For $n$ independent trials with success probability $p^2$
$$
P(a \text { to } b \text { successes }) \approx \Phi\left(\frac{b+\frac{1}{2}-\mu}{\sigma}\right)-\Phi\left(\frac{a-\frac{1}{2}-\mu}{\sigma}\right)
$$
where $\mu=n p$ is the mean, and $\sigma=\sqrt{n p q}$ is the standard deviation.



In physics, the mass-energy equivalence is stated 
by the equation $E=mc^2$, discovered in 1905 by Albert Einstein.

\subsection*{Matricer}
\subsubsection*{bmatrix}
\[\begin{bmatrix}
    3 & 6 & 9\\
    2 & 4 & 6\\
    4 & 5 & 2\\
\end{bmatrix}\]

\subsubsection*{pmatrix}
\[\begin{pmatrix}
    3 & 6 & 9\\
    2 & 4 & 6\\
    4 & 5 & 2\\
\end{pmatrix}\]
\subsubsection*{matrix}
\[\begin{matrix}
    3 & 6 & 9\\
    2 & 4 & 6\\
    4 & 5 & 2\\
\end{matrix}\]
\subsubsection*{Andre typer Matricer:}
https://www.overleaf.com/learn/latex/Matrices




\subsection*{Tuborgparanteser!}
I det følgende kan I se eksempler på, hvad der sker, når man bruger tuborgparenteser! Se eksempelvis når vi forsøger at skrive n opløftet i 21: \\

Uden Tuborgparenteser:
\[n^21\]

Med Tuborgparenteser:
\[n^{21}\]

Her er to eksempler når man skal lave brøker. Når der kun er to tal fungere det ok:

\[\frac{2}{3}\]


\[\frac 2 3\]


Imidlertid opstår der problemer, når vi forsøger med mere komplicerede udtryk! Se følgende:
\[\frac{x*3+5}{4^x}\]

\[\frac x*3+5 4^x\]


\subsection*{Symbol liste}
Symbol Pallet i Overleaf $\Omega$

Eller:

https://tug.ctan.org/info/symbols/comprehensive/symbols-a4.pdf 

Eller:

https://gitlab.gbar.dtu.dk/latex/latex-website-resources/raw/master/guides/latex\_cheatsheet.pdf

\newpage
\subsection{Redskaber til tegn og ligninger}
Det kan være smart at benytte mathpix, så man slipper for at sætte en ligning ind selv.
\subsubsection{Mathpix}



Nedenstående ligning er skrevet vha. MathPix. Oprindeligt var den skrevet i maple:
\begin{equation}
-\frac{1}{4} t^{2}-\frac{1}{2} x+\frac{1}{4} t^{3}+t x
\end{equation}
\begin{equation}
-\frac{1}{4} t^{2}-\frac{1}{2} x+\frac{1}{4} t^{3}+t x
\end{equation}

\begin{figure}[H]
    \centering
    \includegraphics[scale=0.25]{Billeder/mathPix1.png}
    \caption{MathPix på computeren}
    \label{fig:Mathpix_computer}
\end{figure}

\newpage
Man kan også benytte MathPix til at tyde symboler:
\begin{figure}[H]
    \centering
    \includegraphics[scale=0.2]{Billeder/2021-09-09 12-18-39.png}
    \caption{MathPix brugt på telefonen}
    \label{fig:mathpix_telefon}
\end{figure}

\newpage
\subsubsection{Detexify}
For at tyde tegn kan man benytte http://detexify.kirelabs.org/classify.html:
\begin{figure}[H]
    \centering
    \includegraphics[scale=0.4]{Billeder/detextify.png}
    \caption{\url{http://detexify.kirelabs.org/classify.html}, benyttes til at finde latex koden til bl.a. tegn som $\Omega$}
    \label{fig:detextify}
\end{figure}